\clearpage
\section{Text Calibration and Independent Image/Reader Validation}
\label{app:image-readout}

\subsection{Finite calibration of the answer interface}

The pilot motivates two distinct checks: whether answer instructions
preserve known textual meaning, and whether the selected interface and
frozen reader generalize to independent images. The text check uses 16
constructed descriptions: two target objects (red square and blue circle),
four positive/negative patterns, and both truth values. Every prompt begins
with ``The target finding is a \{object\}.'' The object placeholder is
replaced literally throughout. Figure~\ref{fig:text-calibration-prompts}
shows all description patterns and response instructions. Their order and
the routing rule are fixed before observing outcomes.

\begin{figure}[H]
\centering
\fbox{\begin{minipage}{0.94\linewidth}
\small\raggedright
The scene contains a \{object\}. / The scene contains no \{object\}.\par
A \{object\} is present. / A \{object\} is absent.\par
It is false that the \{object\} is absent. / It is false that the \{object\} is present.\par
The list is complete. Listed objects: \{object\}, green triangle. /\par
The list is complete. Listed objects: green triangle.
\end{minipage}}\par
{\small (a) Positive / negative description pairs}\par\medskip
\fbox{\begin{minipage}{0.94\linewidth}
\small\raggedright
Is the finding present?\par
Answer yes or no.
\end{minipage}}\par
{\small (b) Yes/no question and instruction}\par\medskip
\fbox{\begin{minipage}{0.94\linewidth}
\small\raggedright
Answer A if the finding is present and B if it is absent. Reply with A or B only.
\end{minipage}}\par
{\small (c) Incumbent A-present instruction}\par\medskip
\fbox{\begin{minipage}{0.94\linewidth}
\small\raggedright
Choose the correct option.\par
A. The finding is present.\par
B. The finding is absent.\par
Reply with A or B only.
\end{minipage}}\par
{\small (d) Candidate A-present instruction}
\caption{\textbf{Constructed descriptions and answer instructions for text calibration.}
Each description follows the target preface and precedes the question in
(b). The A/B conditions replace only the answer instruction; B-present
swaps the letter assignments in (c), or the option meanings in (d).}
\label{fig:text-calibration-prompts}
\end{figure}

All five conditions share the question ``Is the finding present?'' The
yes/no margin takes the maximum yes singleton logit minus the maximum no
singleton logit; a secondary margin replaces each maximum with
log-sum-exp over the same variants. The A/B margin uses singleton token
IDs 319 and 350 and is oriented toward presence. Each margin must be
strictly positive for presence and strictly negative for absence; zero is
an error. A yes/no route requires both aggregations to be correct on all
16 descriptions and at least one candidate A/B condition to make an error.
This finite constructed-suite rule supplies no population or clinical
inference. The 80 text outcomes take 15 dispatcher seconds on one A100.

\clearpage
Table~\ref{tab:text-calibration} reports every condition. Its minimum
correctly oriented margin multiplies a presence-oriented score by $+1$
for a positive description and by $-1$ for a negative description.
Yes/no passes both aggregations (the secondary minimum is 0.12500056).
Each candidate A/B condition fails the two negative complete-list
descriptions, one per object. The registered rule therefore selects yes/no.

\begin{table}[H]
\centering
\caption{Answer-sign accuracy on constructed descriptions.}
\label{tab:text-calibration}
\resizebox{\ifdim\width>\linewidth\linewidth\else\width\fi}{!}{%
\begin{tabular}{lrrr}
\toprule
Condition & Primary errors / 16 & Minimum oriented margin & LSE errors / 16 \\
\midrule
Yes/no & 0 & 0.125 & 0 \\
Incumbent A-present & 7 & $-3.500$ & --- \\
Incumbent B-present & 8 & $-1.000$ & --- \\
Candidate A-present & 2 & $-2.000$ & --- \\
Candidate B-present & 2 & $-1.500$ & --- \\
\bottomrule
\end{tabular}%
}
\end{table}

Table~\ref{tab:text-pairs} gives the positive-minus-negative primary-margin
change for each description pair, in Figure~\ref{fig:text-calibration-prompts}(a)
order. All changes are positive, including conditions with sign errors.
Sensitivity to a textual truth change therefore does not by itself
establish correct individual answer semantics. These results select an
interface for image testing, not a clinically capable model.

\begin{table}[H]
\centering
\caption{Paired text-margin changes across description patterns.}
\label{tab:text-pairs}
\resizebox{\ifdim\width>\linewidth\linewidth\else\width\fi}{!}{%
\begin{tabular}{llrrrr}
\toprule
Condition & Object & Contains & Present/absent & Negation & Complete list \\
\midrule
Yes/no & Red square & 5.875 & 5.625 & 5.250 & 1.125 \\
 & Blue circle & 5.750 & 5.500 & 6.000 & 1.125 \\
Incumbent A-present & Red square & 3.500 & 3.375 & 3.375 & 0.125 \\
 & Blue circle & 3.375 & 3.125 & 2.875 & 0.250 \\
Incumbent B-present & Red square & 0.750 & 0.250 & 0.875 & 0.250 \\
 & Blue circle & 1.125 & 0.500 & 1.125 & 0.250 \\
Candidate A-present & Red square & 5.500 & 5.875 & 5.750 & 0.250 \\
 & Blue circle & 5.250 & 6.250 & 5.250 & 0.375 \\
Candidate B-present & Red square & 4.250 & 4.500 & 5.375 & 0.375 \\
 & Blue circle & 4.125 & 5.125 & 5.750 & 0.375 \\
\bottomrule
\end{tabular}%
}
\end{table}

\subsection{Independent paired images and frozen measurements}

The image cohort contains validation patients absent from the entire
original 26,229-row manifest. The lexicographically first image identifier
represents each patient. String-sorted patient identifiers are permuted
with NumPy seed 20260916: the first 700 supply index images and the next
700 supply donors, paired by position. All 1,400 patients are distinct and
independent of the pilot cohort. Index labels contain 25 Effusion-positive
and 675 negative cases; donor labels contain 22 positive and 678 negative
cases. Selection is label-blind within this pool. The pilot instead samples
within the original disease-positive/No Finding-sampled manifest, so the
two cohorts do not share the same sampling frame.

Both wordings use the accepted LLaVA-1.5-7B snapshot, native 336-pixel
processor and chat template, configured for BF16 with batch size 16.
The primary score is the first-answer maximum-variant yes-minus-no logit
margin, using singleton yes IDs 3869, 4874, 22483 and no IDs 694, 1939,
11698. Logits are retained in float32; the score's sigmoid is not a
calibrated disease probability. Secondary log-sum-exp scores reuse those
same logits. The run measures 2,800 image outcomes and two fixed-prompt
image-free priors in 260 dispatcher seconds on one A100.

\clearpage
Figure~\ref{fig:image-readout-prompts} gives the exact image prompts.
Let $w\in\{0,1\}$ index the accepted and alternative wording,
respectively. Let $A_w$ be real-image score AUROC against index labels,
$D_w$ donor-image score AUROC against those same index labels, and
$G_w=A_w-D_w$ the image advantage. Donor-own-label AUROC is a separate
descriptive diagnostic, not the donor term in $G_w$.

\begin{figure}[H]
\centering
\fbox{\begin{minipage}{0.94\linewidth}
\small\raggedright
Is there a pleural effusion in this chest radiograph? Answer yes or no.
\end{minipage}}\par
{\small (a) Accepted wording}\par\medskip
\fbox{\begin{minipage}{0.94\linewidth}
\small\raggedright
Does this chest radiograph show a pleural effusion? Answer yes or no.
\end{minipage}}\par
{\small (b) Alternative wording}
\caption{\textbf{Fixed Effusion prompts for independent image validation.}
Each prompt is evaluated with index images, donor images, and once without
an image. The response instruction is shared across wordings.}
\label{fig:image-readout-prompts}
\end{figure}

\subsection{Image opportunity and wording under the primary family}

Table~\ref{tab:image-auroc} presents descriptive 95\% percentile intervals.
The resampling unit is a complete index--donor pair, shared across both
wordings. All 10,000 paired bootstrap draws (seed 20260917) are valid;
AUROC assigns half credit to ties. Each individual interval describes its
own quantity and does not implement the registered joint decision.

\begin{table}[H]
\centering
\caption{Real-image and donor-image answer discrimination.}
\label{tab:image-auroc}
\resizebox{\ifdim\width>\linewidth\linewidth\else\width\fi}{!}{%
\begin{tabular}{lrr}
\toprule
Quantity & Accepted wording & Alternative wording \\
\midrule
Real AUROC $A_w$ & 0.5259 & 0.5089 \\
Descriptive 95\% interval & $[0.4135,0.6367]$ & $[0.3916,0.6255]$ \\
Donor AUROC $D_w$ (index labels) & 0.4072 & 0.3975 \\
Descriptive 95\% interval & $[0.3085,0.5107]$ & $[0.2868,0.5113]$ \\
Image advantage $G_w$ & 0.1187 & 0.1113 \\
Descriptive 95\% interval & $[-0.0013,0.2383]$ & $[-0.0357,0.2583]$ \\
Donor-own-label AUROC & 0.4542 & 0.5117 \\
\bottomrule
\end{tabular}%
}
\end{table}

The registered primary family contains $(A_0-0.5,G_0,G_1-G_0)$.
For each common valid draw, we take the maximum absolute deviation of
these three bootstrap estimates from their observed estimates. The
95th percentile of those maxima, using NumPy's linear quantile, supplies
a common radius of 0.1304. Table~\ref{tab:image-family} reports each point
estimate plus or minus this radius. These are untruncated,
unstudentized simultaneous intervals, not studentized max-$T$ intervals.
The rule requires at least 9,500 draws with all four AUROCs and all three
contrasts finite. Image opportunity requires strictly positive lower
bounds for both $A_0-0.5$ and $G_0$; wording dependence requires the third
interval to exclude zero.

\begin{table}[H]
\centering
\caption{Simultaneous image-opportunity and wording contrasts.}
\label{tab:image-family}
\resizebox{\ifdim\width>\linewidth\linewidth\else\width\fi}{!}{%
\begin{tabular}{lrr}
\toprule
Contrast & Estimate & Simultaneous 95\% interval \\
\midrule
$A_0-0.5$ & 0.0259 & $[-0.1045,0.1563]$ \\
$G_0$ & 0.1187 & $[-0.0117,0.2491]$ \\
$G_1-G_0$ & $-0.0073$ & $[-0.1377,0.1231]$ \\
\bottomrule
\end{tabular}%
}
\end{table}

Neither opportunity lower bound is positive. The cross-wording interval
also spans zero, leaving wording dependence unresolved rather than
establishing no effect or equivalence. The descriptive real-image AUROC
difference $A_1-A_0$ is $-0.0170$ (95\% interval $[-0.0650,0.0300]$).

\clearpage
\subsection{Independent frozen-reader replication}

The Effusion reader uses index images under the accepted wording. Its
features average all 577 tokens at the consumed vision block
\texttt{model.vision\_tower.encoder.layers.22}, stored in float16.
The projection, training mean and scale, clinical coefficients, and all
20 control models and type maps are reused without fitting. Controls
predict fixed random binary labels assigned to recurring
view--sex--age-decade types; each is evaluated against its own type labels,
not against Effusion labels. Selectivity subtracts their mean AUROC from
clinical AUROC, as defined in Section~\ref{sec:protocol}. It is not a direct
estimate of deconfounded clinical information.

Table~\ref{tab:image-reader} compares the pilot and independent cohort.
The frozen-reader analysis uses the same 10,000 bootstrap indices as the
answer analysis but a separate validity mask requiring finite clinical
and all 20 control AUROCs. All draws are valid. The registered replication
criterion requires at least ten cases per class, 9,500 valid draws, and
a positive 2.5th-percentile selectivity bound. That bound is negative, so
the pilot's positive selectivity does not replicate.

\begin{table}[H]
\centering
\caption{Frozen-reader performance across the two cohorts.}
\label{tab:image-reader}
\resizebox{\ifdim\width>\linewidth\linewidth\else\width\fi}{!}{%
\begin{tabular}{lrr}
\toprule
Quantity & Pilot & Independent index cohort \\
\midrule
Effusion-positive / negative & 34 / 666 & 25 / 675 \\
Clinical AUROC & 0.8196 & 0.6664 \\
Mean control AUROC & 0.6758 & 0.6775 \\
Selectivity & 0.1438 & $-0.0111$ \\
Selectivity 95\% interval & $[0.0658,0.2163]$ & $[-0.1541,0.1185]$ \\
Positive lower bound & Yes & No \\
\bottomrule
\end{tabular}%
}
\end{table}

At the point-estimate level, the selectivity difference primarily reflects
lower clinical AUROC, not a higher control mean. The pilot clinical AUROC
interval is $[0.7424,0.8920]$; no between-cohort difference interval was
registered. Different sampling frames, limited positive cases, and
selection of Effusion from the pilot are unresolved explanations, not
identified mechanisms. Reader results neither filter the answer grid nor
alter its simultaneous decision.

Table~\ref{tab:image-controls} lists every frozen control in its stored
order, making the reported mean auditable. Their task-dependent AUROCs
describe recurring nuisance-type predictability and should not be read as
20 alternative clinical Effusion classifiers.

\begin{table}[H]
\centering
\caption{Frozen nuisance-type control AUROCs on independent patients.}
\label{tab:image-controls}
\resizebox{\ifdim\width>\linewidth\linewidth\else\width\fi}{!}{%
\begin{tabular}{rrrrrrrr}
\toprule
Control & AUROC & Control & AUROC & Control & AUROC & Control & AUROC \\
\midrule
1 & 0.9000 & 6 & 0.7416 & 11 & 0.6806 & 16 & 0.7076 \\
2 & 0.5961 & 7 & 0.5952 & 12 & 0.5961 & 17 & 0.5706 \\
3 & 0.6346 & 8 & 0.7585 & 13 & 0.7022 & 18 & 0.6587 \\
4 & 0.7142 & 9 & 0.6728 & 14 & 0.6036 & 19 & 0.7593 \\
5 & 0.7397 & 10 & 0.6238 & 15 & 0.6194 & 20 & 0.6749 \\
\bottomrule
\end{tabular}%
}
\end{table}

\clearpage
\subsection{Descriptive aggregation and prior diagnostics}

Table~\ref{tab:image-aggregation} recomputes discrimination from the same
logits using log-sum-exp (LSE). The accepted-wording image advantage has
a positive descriptive percentile interval, but this secondary endpoint
does not replace the primary simultaneous family. Intervals for the
paired difference from the primary maximum-variant advantage span zero
under both wordings. The evidence does not isolate aggregation choice as
the source of unresolved primary opportunity.

\begin{table}[H]
\centering
\caption{Same-logit aggregation sensitivity.}
\label{tab:image-aggregation}
\resizebox{\ifdim\width>\linewidth\linewidth\else\width\fi}{!}{%
\begin{tabular}{lrr}
\toprule
Quantity & Accepted wording & Alternative wording \\
\midrule
LSE real-image AUROC & 0.5263 & 0.5075 \\
LSE donor AUROC (index labels) & 0.3990 & 0.3944 \\
LSE image advantage & 0.1273 & 0.1131 \\
Descriptive 95\% interval & $[0.0019,0.2504]$ & $[-0.0377,0.2637]$ \\
LSE minus primary advantage & 0.0087 & 0.0017 \\
Descriptive 95\% interval & $[-0.0053,0.0220]$ & $[-0.0137,0.0166]$ \\
\bottomrule
\end{tabular}%
}
\end{table}

Table~\ref{tab:image-priors} reports margin means, Brier error of the
sigmoid score, and total probability mass on the accepted answer tokens.
Real and donor Brier errors both use index labels. The alternative wording
shifts both image-source margins toward yes and has higher Brier error;
its real-image AUROC and paired image advantage remain unresolved under
their respective intervals. This pattern distinguishes a score/prior
shift from established improvement in disease discrimination.

\begin{table}[H]
\centering
\caption{Answer margins, predictive error, and token mass.}
\label{tab:image-priors}
\resizebox{\ifdim\width>\linewidth\linewidth\else\width\fi}{!}{%
\begin{tabular}{llrr}
\toprule
Input & Quantity & Accepted wording & Alternative wording \\
\midrule
Real image & Mean primary margin & 0.0488 & 0.7998 \\
 & Brier error & 0.2728 & 0.4604 \\
 & Mean answer-token mass & 0.9972 & 0.9965 \\
Donor image & Mean primary margin & 0.0380 & 0.7907 \\
 & Brier error & 0.2726 & 0.4596 \\
 & Mean answer-token mass & 0.9972 & 0.9965 \\
No image & Primary margin & 0.8750 & 1.2500 \\
 & Answer-token mass & 0.9932 & 0.9926 \\
 & Constant-score AUROC & 0.5000 & 0.5000 \\
\bottomrule
\end{tabular}%
}
\end{table}

Each image-free prompt produces one constant score, whose AUROC against
index labels is 0.5 under half-credit ties. These are prior and token-mass
diagnostics, not patient-varying no-image discrimination estimates.
Subtracting a wording-specific constant from all its image scores leaves
AUROC unchanged. High answer-token mass establishes concentration on the
chosen verbalizers, not clinical correctness.

Taken together, finite text calibration selects a semantically consistent
yes/no interface for the constructed suite. Independent images establish
neither the registered opportunity conjunction nor wording dependence,
and the frozen clinical reader does not reproduce its pilot selectivity.
The pilot remains an observation within its sampling frame; it does not
support a reproducible reader--answer separation or identify the mechanism
behind the downstream write-control outcomes.

\FloatBarrier
